\documentclass[12pt,a4paper]{article}
\usepackage[utf8]{inputenc}
\usepackage[T5]{fontenc}
\usepackage[vietnamese]{babel}
\usepackage{geometry}
\geometry{a4paper, margin=2cm}
\usepackage{hyperref}



\usepackage{parskip}


\setlength{\parskip}{6pt}

\title{\textbf{Dự án GTS (Go to Success)}}
\author{Phạm Văn Chung -- Khoa học Dữ liệu}

\date{}

\begin{document}

\maketitle



\section*{Lời mở đầu}



Kính thưa Ban Giám khảo,


Em là Phạm Văn Chung, sinh ra ở một vùng quê nghèo Yên Bái. Nhà em không có máy tính, trường em cũng không có phòng tin học. Lên lớp 10, em mới biết đến ``lập trình'' qua một cuốn sách mượn từ thư viện.


Không có máy tính, em viết code ra vở ô ly, tưởng tượng nó chạy trong đầu. Có hôm em thức dậy lúc 4 giờ sáng ra quán internet để gõ thử những dòng mình đã viết. Một lỗi quên dấu chấm phẩy từng tốn của em cả buổi chiều và mấy chục nghìn -- cả tuần tiền tiêu vặt thời đó.


Nhưng chính những ngày cô đơn ấy đã nhen nhóm trong em một ý tưởng: \textit{Tại sao không tạo ra một người thầy ảo luôn sẵn sàng giúp đỡ, miễn phí, cho tất cả học sinh nghèo như mình?}


\section*{Hành trình GTS}


Cuối lớp 12, em bắt tay xây dựng GTS một mình -- trên chiếc laptop cũ mượn của anh họ. Dự án đạt \textbf{Giải Nhì Khoa học Kỹ thuật cấp Tỉnh}. Giải thưởng ấy khẳng định: một học sinh vùng cao, chưa từng qua lớp học nào, vẫn có thể tạo ra sản phẩm công nghệ có giá trị.


Hôm nay, là sinh viên ngành Khoa học Dữ liệu, em tiếp tục phát triển GTS -- tích hợp AI, xây dựng bảo mật, và giữ vững cam kết: \textbf{miễn phí mãi mãi}.


\section*{GTS giải quyết ba vấn đề}



























\begin{enumerate}





  \item \textbf{Rào cản thiết bị:} GTS chạy mượt trên cả điện thoại cũ, máy tính cấu hình thấp, mạng chậm.
  \item \textbf{Không người hướng dẫn:} \textit{AI Mentor} -- người thầy 24/7, chỉ dẫn cách nghĩ, không giải bài hộ.
  \item \textbf{Rào cản tài chính:} Hoàn toàn miễn phí. Mọi học sinh, dù ở đâu, đều có quyền tiếp cận như nhau.
\end{enumerate}


\section*{GTS là gì?}


Một nền tảng học lập trình trực tuyến, nơi học sinh viết code, nộp bài, chấm điểm tự động -- như một kỳ thi Olympic thực thụ. Hệ thống gồm: giao diện Web (JavaScript/Firebase), máy chủ chấm bài (Python/Flask), AI Mentor (DeepSeek/Grok), và cơ sở dữ liệu thời gian thực (Firebase).








Chỉ cần một chiếc điện thoại có internet, học sinh ở bất kỳ đâu cũng có thể lập trình, nhận kết quả trong vài giây và được AI hướng dẫn tận tình.


\section*{Tầm nhìn}













Em không tham vọng xây kỳ lân công nghệ. Em chỉ muốn \textbf{không còn bạn nào phải một mình vật lộn với lỗi code, bỏ cuộc giữa chừng vì không ai chỉ dạy}.


Những tin nhắn các em học sinh gửi về -- \textit{``Nhờ GTS em mới biết lập trình''}, \textit{``Em ở Hà Giang, AI Mentor như người bạn đồng hành''} -- là động lực để em tiếp tục.


\section*{Lời kết}


Em viết những dòng này không để khoe về sản phẩm hoàn hảo. Em viết để kể câu chuyện của một cậu bé vùng cao không có gì, nhưng có niềm tin: \textbf{công nghệ có thể thay đổi cuộc đời, và ai cũng xứng đáng có cơ hội}.







Học bổng Bản lĩnh thế hệ tôn vinh tinh thần kiên cường và khát vọng tạo ảnh hưởng tích cực. Em mong Ban Giám khảo trao cơ hội để em tiếp tục hành trình đưa công nghệ đến với mọi học sinh Việt Nam, bất kể các em ở đâu, hoàn cảnh ra sao.









































Trân trọng cảm ơn!

\textbf{Phạm Văn Chung}

\end{document}
